% Generated draft: requires mathematical and source review.
\documentclass[11pt]{article}
\usepackage[margin=25mm]{geometry}
\usepackage{fontspec}
\setmainfont{Latin Modern Roman}
\usepackage{amsmath,amssymb,mathtools,unicode-math}
\setmathfont{Latin Modern Math}
\usepackage{enumitem,needspace,xurl,hyperref}
\setlength{\parindent}{0pt}
\setlength{\parskip}{5pt}
\setlength{\emergencystretch}{3em}
\newcommand{\N}{\mathbb{N}}
\newcommand{\Z}{\mathbb{Z}}
\newcommand{\Q}{\mathbb{Q}}
\newcommand{\R}{\mathbb{R}}
\newcommand{\C}{\mathbb{C}}
\newcommand{\F}{\mathbb{F}}
\newcommand{\A}{\mathbb{A}}
\newcommand{\PP}{\mathbb{P}}
\newcommand{\E}{\mathbb{E}}
\newcommand{\eps}{\varepsilon}
\newcommand{\dd}{\,\mathrm{d}}
\begin{document}
\section*{UnsolvedMath ID: 30001912}
\textbf{Source problem code:} catalogue entry
\subsection{Definitions and mathematical statement}
For a field \(k\) and a specified number \(n\) of variables, a Laurent polynomial is an element
\[
 f\in k[x_1^{\pm1},\ldots,x_n^{\pm1}],\qquad
 f=\sum_{m\in A}c_m x_1^{m_1}\cdots x_n^{m_n},
\]
where \(A\subset\Z^n\) is finite and \(c_m\in k\). Under the usual algebraic convention, a Fano manifold is a smooth projective geometrically integral variety \(X\) whose anticanonical line bundle \(\mathcal O_X(-K_X)\) is ample. The source does not specify the base field, the number of variables, the dimension of \(X\), or any additional restrictions on \(f\) or \(X\); these conventions therefore do not resolve the intended geometric setting.

The source question is: \emph{What necessary conditions must a Laurent polynomial \(f\) satisfy in order to be mirror to a Fano manifold?}

Once a geometric setting and a precise mirror relation \(\mathcal M(f,X)\) have been specified, the requested conditions can be expressed as properties \(P\) of Laurent polynomials for which
\[
 \forall f\in k[x_1^{\pm1},\ldots,x_n^{\pm1}],\qquad
 \bigl(\exists X\text{ a Fano manifold with }\mathcal M(f,X)\bigr)
 \Longrightarrow P(f).
\]
Here \(\mathcal M\) is only a placeholder for the unspecified meaning of being mirror; no definition of that relation is supplied or asserted. The class of admissible \(X\), including any requirement that \(\dim X=n\), must also be specified before this schema becomes a determinate mathematical problem. The question asks for necessary conditions, without asserting their sufficiency or asking explicitly for a complete characterization.
\subsection{Short English statement}
Determine restrictions that every Laurent polynomial mirror to a Fano manifold must satisfy. The supplied question leaves the meaning of mirror and the geometric setting unspecified, so a precise version requires those choices. Its supplied status is partially solved; this is a dated curation label, not an independently verified conclusion.
\subsection{Sources}
\begin{itemize}
\item Local UnsolvedMath registry, permanent ID \(30001912\), source problem code \texttt{catalogue entry}, title \emph{Necessary Conditions for Laurent-Polynomial Fano Mirrors}, in the Oberwolfach Reports Open Problems collection. The record gives proposed year 2011 and names no individual proposer.
\item Supplied report citation: \emph{Complex Algebraic Geometry (2012)}, Oberwolfach Reports, DOI 10.4171/owr/2011/47, pp.\ 2669--2733. \url{https://doi.org/10.4171/owr/2011/47}. The background describes these reports as documenting mathematical workshops and their open-problem discussions. The supplied citation uses 2012, whereas the proposed year and DOI contain 2011; this discrepancy is retained without resolution.
\item Supplied literature reference: Coates--Corti--Galkin--Kasprzyk, quantum periods. \url{https://arxiv.org/abs/1303.3288}.
\item Supplied literature reference: Coates--Kasprzyk--Pitton--Tveiten, maximally mutable Laurent polynomials. \url{https://arxiv.org/abs/2107.14253}.
\item Supplied literature reference: original OWR report. \url{https://doi.org/10.4171/owr/2011/47}.
\item Supplied literature assessment, checked at \texttt{2026-08-21}: the status is partially solved. The assessment states that period, regularity, and extremality conditions and large families of Laurent-polynomial mirrors are established, but that there is no complete necessary-and-sufficient characterization of mirrors to Fano manifolds. This is supplied dated curation derived from the OWR record, not independently verified here; its time-sensitive status and cited claims require independent checking.
\item The registry labels the statement \texttt{corrected\_verified} and reports that its mathematical content was checked against the official report, with notation, context, and question form normalized without changing the source-level claim. This is supplied verification metadata, not a claim that the report was consulted for this draft.
\end{itemize}
\end{document}
